\documentclass[a4paper]{article}     
\usepackage{amsfonts}
\usepackage{amsmath}
\usepackage{amssymb}
\usepackage{graphicx}
\usepackage{color}

\newcommand{\Q}{\mathbb{Q}}
\newcommand{\R}{\mathbb{R}}
\newcommand{\llim}{\lim\limits}
\newcommand{\dint}{\displaystyle\int}

\begin{document}

\noindent \large Auteur: Ilyas Sefiani \\
\noindent \large Studierichting: Informatica\\
\noindent \large Oefeningenreeks 4A nr. 48.3\\

\medskip

\normalsize

$\displaystyle\int \ln(x^2-1)dx$\\

$\left \{
\begin{array}{l}
u=\ln(x^2-1) \\
dv=dx
\end{array}
\right .
\Rightarrow
\left \{
\begin{array}{l}
du=\dfrac{2x}{x^2-1}dx \\
v=x
\end{array}
\right .$\\

$=x\ln(x^2-1) - \displaystyle\int \dfrac{2x^2}{x^2-1}dx$\\

$=x\ln(x^2-1) - 2\displaystyle\int \dfrac{x^2-1+1}{x^2-1}dx$\\

$=x\ln(x^2-1) - 2\left(\displaystyle\int dx + \displaystyle\int \dfrac{dx}{x^2-1}\right)$\\

$=x\ln(x^2-1) - 2\left(x + \displaystyle\int \dfrac{dx}{x^2-1}\right)$\\

Fuss toepassen:\\

$\displaystyle\int \dfrac{dx}{x^2-1} = \dfrac{A}{x-1} + \dfrac{B}{x+1}$\\

Voor $x = 1$:\\

$A = \dfrac{1}{1+1} = \dfrac{1}{2}$\\

Voor $x = -1$:\\

$B = \dfrac{1}{-1-1} = -\dfrac{1}{2}$\\

$=x\ln(x^2-1) - 2\left(x + \dfrac{1}{2}\displaystyle\int \dfrac{dx}{x-1}-\dfrac{1}{2}\displaystyle\int \dfrac{dx}{x+1}\right)$\\

$=x\ln(x^2-1) - 2\left(x + \dfrac{1}{2}\ln|x-1|-\dfrac{1}{2}\ln|x+1|\right) + C$\\

$=x\ln(x^2-1) - 2x -\ln|x-1|+\ln|x+1| + C$\\

$=x\ln(x^2-1) - 2x +\ln\left|\dfrac{x+1}{x-1}\right| + C$\\

\begin{thebibliography}{999}
\end{thebibliography}
\end{document} 